\documentclass[12pt]{article}
\usepackage{lmodern}
\usepackage{geometry}
\usepackage{enumitem}
\setlist[description]{leftmargin=1cm,labelindent=1cm}
\usepackage{amsmath, amsthm, amssymb, mathtools}
\usepackage[hidelinks]{hyperref}
\usepackage[portuguese]{babel}

\theoremstyle{plain}
\newtheorem{thm}{Teorema}[section]
\newtheorem{prop}[thm]{Proposição}
\newtheorem{lem}[thm]{Lema}
\newtheorem{cor}[thm]{Corolário}

\theoremstyle{definition}
\newtheorem{defn}[thm]{Definição}
\newtheorem{exmp}[thm]{Exemplo}
\newtheorem{prblm}[thm]{Problema}
\newtheorem{prcp}[thm]{Princípio}

\theoremstyle{remark}
\newtheorem*{rem}{Nota}

\geometry{
 a4paper,
 left=3cm,
 top=2cm,
 right=2cm,
 bottom=2cm
}

\begin{document}
\begin{center}
    \textbf{\large UNIVERSIDADE FEDERAL DO RIO GRANDE DO SUL} \\
    \textbf{\large INSTITUTO DE MATEMÁTICA E ESTATÍSTICA} \\
    \vspace{3cm}
    \textbf{\Large MÉTODOS APLICADOS II} \\
    \vspace{1cm}
    \textbf{\large LISTA DE RESULTADOS} \\
    \vspace{3cm}

    \textbf{LEONARDO SUAREZ}

    \vfill
    \textbf{\large Porto Alegre, RS} \\
    \textbf{\large \today}
\end{center}

\newpage
\tableofcontents
\newpage

\section{Soluções por Séries de Potências}

\begin{rem} (fatos importantes sobre séries de potências)
    \begin{itemize}
        \item \textbf{Definição.} Uma série de potências centrada em \(x-a\) é uma série infinita da forma \begin{equation}\sum_{n=0}^{\infty} c_n (x-a)^n = c_0 + c_1(x-a) + c_2(x-a)^2 + \cdots\end{equation}
        \item \textbf{Convergência.} Uma série de potências é convergente no ponto \(x=x_0\) se a sua sequência de somas parciais \(\{S_N(x)\}\) converge para um número real para \(N\to\infty\), ou seja, \(\displaystyle\lim_{N\to\infty} S_N(x_0) = \lim_{N\to\infty} \sum_{n=0}^{N} c_n (x-a)^n\) existe. Se o limite não existe, a série diverge em \(x_0\).
        \item \textbf{Intervalo de convergência.} O conjunto de todos os pontos \(x\) para os quais a série de potências converge é chamado de intervalo de convergência. O centro do intervalo de convergência é o centro \(a\) da série de potências.
        \item \textbf{Raio de convergência.} O raio \(R\) do intervalo de convergência de uma série de potências é chamado de raio de convergência. Se a série converge para todo \(x\), então \(R=\infty\). Se \(R>\), então a série converge para \(|x-a|<R\) e diverge para \(|x-a|>R\). Se \(R=0\), então a série converge somente em \(x=a\). Uma série de potências pode ou não convergir nos pontos finais do intervalo de convergência.
        \item \textbf{Convergência absoluta.} Dentro de seu intervalo de convergência, uma série de potências converge absolutamente, ou seja, a série \(\displaystyle\sum_{n=0}^{\infty} |c_n (x-a)^n|\) também converge para todo \(x \in (a-R, a+R)\).
        \item \textbf{Teste da razão.} A convergência de uma série de potências pode ser determinada pelo teste da razão. Suponha que \(c_n\neq0\) para todo \(n\) em \(\sum_{n=0}^{\infty} c_n (x-a)^n\), e que 
        \begin{equation}
            \lim_{n\to\infty} \left| \frac{c_{n+1}(x-a)^{n+1}}{c_n (x-a)^n} \right| = |x-a| \lim_{n\to\infty} \left| \frac{c_{n+1}}{c_n} \right| = L.
        \end{equation}
        \noindent Se \(L<1\), a série converge absolutamente; se \(L>1\), a série diverge; se \(L=1\), o teste é inconclusivo. O teste da razão é sempre inconclusivo nos pontos \(x=a\pm R\).
        \item \textbf{Uma série de potências define uma função.} Uma série de potências \(\displaystyle\sum_{n=0}^{\infty} c_n (x-a)^n\) define uma função \(f\) por \(\displaystyle f(x) = \sum_{n=0}^{\infty} c_n (x-a)^n\) para todo \(x\) no intervalo de convergência. Se o raio de convergência é \(R>0\) ou \(R = \infty\), então \(f\) é contínua, diferenciável e integrável nos intervalos \((a-R, a+R)\) ou \((-\infty, \infty)\), respectivamente. Além disso, \(f\) pode ser diferenciada e integrada termo a termo. Convergência em um ponto final do intervalo pode ser perdida após a diferenciação ou integração. Se 
        \[ f(x) = \sum_{n=0}^{\infty} c_n x^n = c_0 + c_1 x + c_2 x^2 + \cdots \]
        \noindent é uma série de potências em \(x\), então suas primeiras duas derivadas são \(\displaystyle y' = \sum_{n=1}^{\infty} n c_n x^{n-1}\) e \(\displaystyle y'' = \sum_{n=2}^{\infty} n(n-1)c_n x^{n-2}\).
        \item \textbf{Identidade.} Se \(\displaystyle\sum_{n=0}^{\infty} c_n (x-a)^n = 0\) para todo \(x\) em um intervalo, então \(c_n = 0\) para todo \(n\).
        \item \textbf{Analítica em um ponto.} Uma função \(f\) é dita analítica em um ponto \(x=a\) se existe uma série de potências centrada em \(x-a\) que converge para \(f(x)\) em algum intervalo \((a-R, a+R)\), onde \(R>0\).
    \end{itemize}
\end{rem}

\begin{defn}
    Um ponto \(x=x_0\) é chamado de \emph{ponto ordinário} da equação diferencial
    \begin{equation}\label{eq:edo1}
        a_2(x)y'' + a_1(x)y' + a_0(x)y = 0,
    \end{equation}
    \noindent se ambos os coeficientes \(P(x)\) e \(Q(x)\) na forma padrão
    \begin{equation}\label{eq:edo2}
        y'' + P(x)y' + Q(x)y = 0,
    \end{equation}
\noindent são analíticos em \(x_0\). Caso contrário, \(x_0\) é chamado de \emph{ponto singular}.
\end{defn}

\begin{thm} (existência de soluções por séries de potências)
    Se \(x_0\) é ponto ordinário da equação \eqref{eq:edo1}, então podemos encontrar duas soluções linearmente independentes na forma de séries de potências centradas em \(x_0\), ou seja,
    \[
    y = \sum_{n=0}^{\infty} c_n (x - x_0)^n.
    \]
    Uma solução de série de potência converge pelo menos no intervalo \(|x - x_0| < R\), onde \(R\) é a distância do ponto \(x_0\) ao ponto singular mais próximo, ou também a \emph{cota inferior} do raio de convergência.
\end{thm}

\begin{defn}
    Uma ponto singular \(x=x_0\) é chamado de \emph{ponto singular regular} da equação \eqref{eq:edo1} se as funções \(p(x)=(x-x_0)P(x)\) e \(q(x)=(x-x_0)^2Q(x)\) são analíticas em \(x_0\). Caso contrário, \(x_0\) é chamado de \emph{ponto singular irregular}.
\end{defn}

\begin{thm} (teorema de Frobenius)
    Se \(x_0\) é ponto singular regular da equação \eqref{eq:edo1}, então existe pelo menos uma solução da forma
    \begin{equation}
        y = (x - x_0)^r \sum_{n=0}^{\infty} c_n (x - x_0)^n = \sum_{n=0}^{\infty} c_n (x - x_0)^{n+r},
    \end{equation}
    \noindent onde \(r\) é a raiz de uma equação indicial associada. A série converge pelo menos em algum intervalo \(0<x-x_0<R\).
\end{thm}

\begin{defn}
    Assuma que \(\nu\geq0\). Então,
    \begin{equation}
        x^2y'' + xy' + (x^2 - \nu^2)y = 0
    \end{equation}
    \noindent é chamada de \emph{equação de Bessel de ordem \(\nu\)}.
\end{defn}

\begin{defn}
    Considerando \(n\) um inteiro não-negativo, temos 
    \begin{equation}
        (1-x^2)y'' - 2xy' + n(n+1)y = 0
    \end{equation}
    \noindent que é chamada de \emph{equação de Legendre de ordem \(n\)}.
\end{defn}

\section{Séries de Fourier}

\begin{defn} 
    O \emph{produto interno} de duas funções \(f_1\) e \(f_2\) em um intervalo \([a,b]\) é o número
    \begin{equation}
        \langle f_1, f_2 \rangle = \int_a^b f_1(x)f_2(x) \, dx.
    \end{equation}
\end{defn}

\begin{defn}
    Duas funções \(f_1\) e \(f_2\) são ditas \emph{ortogonais} em \([a,b]\) se
    \begin{equation}
        \langle f_1, f_2 \rangle = 0.
    \end{equation}
\end{defn}

\begin{defn}
    Um conjunto de funções reais \(\{\phi_1(x), \phi_2(x), \ldots\}\) é \emph{ortogonal} em \([a,b]\) se
    \begin{equation}
        \langle \phi_m, \phi_n \rangle = \int_a^b \phi_m(x)\phi_n(x) \, dx = 0, \quad m \neq n.
    \end{equation}
\end{defn}

\begin{defn}
    A \emph{norma quadrada} de uma função \(\phi_n\) é \(\Vert \phi _n \Vert^2 = \langle \phi_n, \phi_n \rangle = \int_a^b \phi_n(x)^2 \, dx\), e a \emph{norma} é dada por \(\Vert \phi_n \Vert = \sqrt{\langle \phi_n, \phi_n \rangle} = \sqrt{\int_a^b \phi_n(x)^2 \, dx}\). Se \(\Vert \phi_n \Vert = 1\), então \(\{\phi_n\}\) é um conjunto \emph{ortonormal} no intervalo \([a,b]\).
\end{defn}

\begin{defn}
    Um conjunto de funções reais \(\{\phi_0(x), \phi_1(x), \phi_2(x), \ldots\}\) é \emph{ortogonal com respeito a uma função peso} \(w(x)\) em \([a,b]\) se
    \begin{equation}
        \int_{a}^{b} w(x)\phi_m(x)\phi_n(x) \, dx = 0, \quad m \neq n.
    \end{equation}

    \noindent A série 
    \begin{equation}
    f(x)=\sum_{n=0}^{\infty} c_n \phi_n(x)\,, \quad c_n=\frac{\int_{a}^{b} f(x)w(x)\phi_n(x) \, dx}{\Vert \phi_n(x)\Vert^2}.
    \end{equation}
    \noindent é chamada de \emph{série de Fourier generalizada} de \(f\).
\end{defn}

\begin{defn}
    A \emph{série de Fourier} de uma função \(f\) definida no intervalo \((-p,p)\) é dada por
    \begin{equation}
        f(x) = \frac{1}{2}a_0 + \sum_{n=1}^{\infty} \left( a_n \cos\frac{n\pi x}{p} + b_n \sin\frac{n\pi x}{p} \right),
    \end{equation}
    onde os coeficientes \(a_n\) e \(b_n\) são dados por
    \begin{align}
        a_0 &= \frac{1}{p} \int_{-p}^{p} f(x) \, dx, \\
        a_n &= \frac{1}{p} \int_{-p}^{p} f(x) \cos\frac{n\pi x}{p} \, dx, \\
        b_n &= \frac{1}{p} \int_{-p}^{p} f(x) \sin\frac{n\pi x}{p} \, dx.
    \end{align}
\end{defn}

\begin{defn}
    Uma função \(f\) é dita \emph{contínua por partes} em um intervalo \([a,b]\) se existe um número finito de pontos de descontinuidade \(x_1<x_2<\cdots<x_n\) em \([a,b]\) e \(f\) é contínua em cada subintervalo \((x_k, x_{k+1})\) para \(k=0,1,\ldots,n\). Consequentemente,
    \begin{itemize}
        \item Os limites laterais \(f(x+)\) e \(f(x-)\) existem para todo \(x \in (a,b)\).
        \item Os limites laterais \(f(a+)\) e \(f(b-)\) existem, mas não é necessário que \(f\) seja contínua ou mesmo definida em \(a\) ou em \(b\)
    \end{itemize}
\end{defn}

\begin{thm} (condições para convergência)
    Sejam \(f\) e \(f'\) contínuas por partes em \([-p,p]\). Então para todo \(x \in (-p,p)\), a série de Fourier de \(f\) converge para \(f(x)\) em um ponto de continuidade. Em um ponto de descontinuidade, a série de Fourier converge para a média
    \begin{equation}
        \frac{f(x+) + f(x-)}{2},
    \end{equation}
    \noindent onde \(f(x+)\) e \(f(x-)\) são os limites laterais de \(f\) em \(x\), respectivamente.
\end{thm}

\begin{thm} (propriedades de funções pares/ímpares)
    \begin{enumerate}
        \item[(a)] O produto de duas funções pares é par.
        \item[(b)] O produto de duas funções ímpares é par.
        \item[(c)] O produto de uma função par por uma função ímpar é ímpar.
        \item[(d)] A soma (diferença) de duas funções pares é par.
        \item[(e)] A soma (diferença) de duas funções ímpares é ímpar.
        \item[(f)] Se \(f\) é par, então \(\int_{-a}^{a} f(x) \, dx = 2\int_{0}^{a} f(x) \, dx\).
        \item[(g)] Se \(f\) é ímpar, então \(\int_{-a}^{a} f(x) \, dx = 0\).
    \end{enumerate}
\end{thm}

\begin{defn} (séries de Fourier para funções pares/ímpares)
    \begin{enumerate}
        \item[(i)] A série de Fourier de uma função par \(f\) definida no intervalo \((-p,p)\) é a \textbf{série em cossenos}
        \begin{equation}
            f(x) = \frac{a_0}{2} + \sum_{n=1}^{\infty} a_n \cos\left(\frac{n\pi x}{p}\right),
        \end{equation}
        \noindent onde 
        \begin{align}
            a_0 &= \frac{2}{p} \int_{0}^{p} f(x) \, dx, \\
            a_n &= \frac{2}{p} \int_{0}^{p} f(x) \cos\left(\frac{n\pi x}{p}\right) \, dx.
        \end{align}
        \item[(ii)] A série de Fourier de uma função ímpar \(f\) definida no intervalo \((-p,p)\) é a \textbf{série em senos}
        \begin{equation}
            f(x) = \sum_{n=1}^{\infty} b_n \sin\left(\frac{n\pi x}{p}\right),
        \end{equation}
        \noindent onde
        \begin{equation}
            b_n = \frac{2}{p} \int_{0}^{p} f(x) \sin\left(\frac{n\pi x}{p}\right) \, dx.
        \end{equation}
    \end{enumerate}
\end{defn}

\renewcommand{\arraystretch}{1.15}
\begin{rem} (revisão de EDOs) \newline
    \begin{tabular}[c]{ll}
        \textbf{Equação de coeficientes constantes} & \textbf{Soluções gerais} \\
        \hline
        \(y'+\alpha y=0\) & \(y=ce^{-\alpha x}\) \\
        \(y''+\alpha^2 y=0,\quad \alpha>0\) & \(y=c_1\cos(\alpha x)+c_2\sin(\alpha x)\) \\
        \(y''-\alpha^2 y=0,\quad \alpha>0\) & \(\begin{cases} y=c_1e^{-\alpha x}+c_2e^{\alpha x} \\ y=c_1\cosh(\alpha x)+c_2\sinh(\alpha x) \end{cases}\) \\
        \\ \textbf{Equação de Cauchy-Euler} & \textbf{Soluções gerais, \(x>0\)} \\
        \hline
        \(x^2y''+ xy'- \alpha^2 y=0,\quad \alpha\geq 0\) & \(\begin{cases} y=c_1x^{-\alpha}+c_2x^{\alpha},\quad\alpha>0 \\ y=c_1+c_2\ln(x),\quad \alpha=0 \end{cases}\) \\
        \\ \textbf{Equação paramétrica de Bessel \(\nu>0\)} & \textbf{Solução geral, \(x>0\)} \\
        \hline
        \(x^2y''+ xy' + \alpha^2 x^2 y=0\) & \(y=c_1J_0(\alpha x)+c_2Y_0(\alpha x)\) \\
        \\ \textbf{Equação de Legendre \(n=0,1,2,\ldots\)} & \textbf{Soluções particulares são polinômios} \\
        \hline
        \((1-x^2)y'' - 2xy' + n(n+1)y=0\) & \(\begin{cases}
            y=P_0(x)=1,\\
            y=P_1(x)=x,\\
            y=P_2(x)=\frac{1}{2}(3x^2-1),\ldots
        \end{cases}\)
    \end{tabular}
\end{rem}

\begin{defn}
    Sejam \(p\), \(q\), \(r\) e \(r'\) funções reais contínuas em um intervalo \([a,b]\), e seja \(r(x)>0\) e \(p(x)>0\) para todo \(x \in [a,b]\). O problema de valor de contorno
    \begin{equation}
        \frac{d}{dx}\left[p(x)\frac{dy}{dx}\right] + q(x)y + \lambda r(x)y = 0, \quad a < x < b,
    \end{equation}
    \noindent sujeito às condições de contorno lineares e homogêneas
    \begin{align}
        A_1 y(a) + A_2 y'(a) &= 0, \\
        B_1 y(b) + B_2 y'(b) &= 0,
    \end{align}
    \noindent onde \(A_1\), \(A_2\), \(B_1\) e \(B_2\) são constantes tais que \(A_1^2 + A_2^2 \neq 0\) e \(B_1^2 + B_2^2 \neq 0\), é chamado de \emph{problema de Sturm-Liouville regular}. Os valores de \(\lambda\) para os quais existem soluções não-triviais são chamados de \emph{autovalores}, e as soluções correspondentes são chamadas de \emph{autovetores}.
\end{defn}

\begin{thm} (propriedades do problema de Sturm-Liouville regular)
    \begin{enumerate}
        \item[(a)] Existe um número infinito de autovalores reais que podem ser ordenados em uma sequência crescente \( \lambda_1 < \lambda_2 < \lambda_3 < \ldots \) tal que \( \lim_{n \to \infty} \lambda_n = +\infty \).
        \item[(b)] A cada autovalor \( \lambda_n \) corresponde um único autovetor \( y_n(x) \), a menos de um fator multiplicativo constante.
        \item[(c)] Autovetores correspondentes a autovalores distintos são linearmente independentes.
        \item[(d)] O conjunto de autovetores correspondente ao conjunto de autovalores é ortogonal com respeito à função peso \( p(x) \) no intervalo \([a,b]\).
    \end{enumerate}
\end{thm}

\section{Problemas de Valor de Contorno em Coordenadas Retangulares}
\section{Problemas de Valor de Contorno em Outros Sistemas de Coordenadas}
\section{Transformadas Integrais}

\begin{defn}
    A \emph{integral de Fourier} de uma função \(f\) definida no intervalo \(\left(-\infty, \infty\right)\) é dada por
    \begin{equation}
        f(x) = \frac{1}{\pi} \int_{0}^{\infty} \left[A(\alpha) \cos(\alpha x) + B(\alpha) \sin(\alpha x)\right] d\alpha,
    \end{equation}
    \noindent onde
    \begin{align}
        A(\alpha) &= \int_{-\infty}^{\infty} f(x) \cos(\alpha x) \, dx, \\
        B(\alpha) &= \int_{-\infty}^{\infty} f(x) \sin(\alpha x) \, dx.
    \end{align}
\end{defn}

\begin{thm} (condições para convergência da integral de Fourier)
    Seja \(f\) e \(f'\) contínuas por partes em todo intervalo finito e seja \(f\) absolutamente integrável em \((-\infty, \infty)\). Então, a integral de Fourier de \(f\) converge para \(f(x)\) em um ponto de continuidade. Em um ponto de descontinuidade, a integral de Fourier converge para a média
    \begin{equation}
        \frac{f(x+) + f(x-)}{2},
    \end{equation}
    \noindent onde \(f(x+)\) e \(f(x-)\) são os limites laterais de \(f\) em \(x\), respectivamente.
\end{thm}

\begin{defn} (integrais de Fourier em senos e cossenos)
    \begin{enumerate}
        \item[(i)] A integral de Fourier de uma função par \(f\) definida no intervalo \(\left(-\infty, \infty\right)\) é a \textbf{integral em cossenos}
        \begin{equation}
            f(x) = \frac{2}{\pi} \int_{0}^{\infty} A(\alpha) \cos(\alpha x) \, d\alpha,
        \end{equation}
        \noindent onde
        \begin{align}
            A(\alpha) &= \int_{0}^{\infty} f(x) \cos(\alpha x) \, dx.
        \end{align}
        \item[(ii)] A integral de Fourier de uma função ímpar \(f\) definida no intervalo \(\left(-\infty, \infty\right)\) é a \textbf{integral em senos}
        \begin{equation}
            f(x) = \frac{2}{\pi} \int_{0}^{\infty} B(\alpha) \sin(\alpha x) \, d\alpha,
        \end{equation}
        \noindent onde
        \begin{align}
            B(\alpha) &= \int_{0}^{\infty} f(x) \sin(\alpha x) \, dx.
        \end{align}
    \end{enumerate}
\end{defn}

\end{document}